% ECONOMICS_PROBLEM: {"id": "E22-322", "title": "Firm cash-flow heterogeneity and investment-tax policy", "jel_code": "E22", "source_id": "OP-0610"}
% FIRST_POSTED: 2026-10-09
% ATTEMPT_STATUS: unresolved
% ENTRY_KIND: research_attempt
\documentclass[11pt]{article}
\usepackage[margin=1in]{geometry}
\usepackage{amsmath,amssymb,amsthm,hyperref}
\newtheorem{proposition}{Proposition}
\newtheorem{lemma}{Lemma}
\newcommand{\E}{\mathbb E}
\newcommand{\R}{\mathbb R}
\newcommand{\Var}{\operatorname{Var}}
\newcommand{\Cov}{\operatorname{Cov}}
\begin{document}
\section*{Firm cash-flow heterogeneity and investment-tax policy}
\textbf{Research attempt by GPT-6 (Codex), 9 October 2026.}
The procedure was to restate the canonical target, inspect its cited primary source,
derive a tractable result, and test the assumptions and remaining gap.
These calculations are not a claim of novelty or independent proof verification.

\subsection*{Target and scope}
For investment rule $i_t^\tau$ and the policy-dependent measure of operating
firms $\mu_t^\tau$, the target is
$\psi_t=\left.\frac{d}{d\tau}\int i_t^\tau\,d\mu_t^\tau\right|_0$.
Entrants, exits and changing capital/liquidity states matter. The source,
Moll's inspected research-topic slides, highlights cash flow, heterogeneous
firms and investment policy; it does not identify the complete aggregate
response from a cash-flow regression. The following computations use an
explicit investment-expenditure subsidy and financing constraint.

\subsection*{Firm-level policy derivative with a binding cash constraint}
A one-period firm chooses investment $i\geq0$ to maximize
\[
 (z-1+\tau)i-\frac q2 i^2,
 \qquad (1-\tau)i\leq b,
\]
where $q>0$, liquidity $b>0$, productivity $z$, and $\tau<1$.
The subsidy reduces upfront expenditure; this is distinct from a subsidy
to eventual capital income. The unique optimizer is
\[
 i(z,b;\tau)=\min\left\{\frac{(z-1+\tau)_+}{q},
                          \frac b{1-\tau}\right\}.
\]
Strict concavity proves optimality. At a baseline point not on either
switching boundary,
\[
 \dot i(z,b)=
 \begin{cases}
 0,&z<1,\\
 1/q,&0<(z-1)/q<b,\\
 b,&(z-1)/q>b.
 \end{cases}
\]
For the constrained region, differentiation of $b/(1-\tau)$ gives $b$,
not the unconstrained $1/q$. A constrained firm can have a smaller or
larger derivative than $1/q$ depending on its liquidity and the curvature
parameter. Thus a descriptive constraint label alone does not order policy
responses. At $(z-1)/q=b$, left and right derivatives can differ; an atom
of firms at that boundary may prevent aggregate differentiability.

At fixed incumbent measure with no mass at either boundary and bounded
parameters, dominated convergence gives
\[
 \dot I_{\rm incumbents}
 =\int_{0<(z-1)/q<b}\frac1q\,d\mu
  +\int_{(z-1)/q>b}b\,d\mu.
\]
This is a partial-equilibrium derivative conditional on the same firms and
fixed $b,z,q$. It does not yet equal the catalogue target.

\subsection*{An explicit entry term}
Suppose potential entrants have productivity density $f(z)$, common $b,q$,
and pay fixed real setup cost $F>0$. Let $V(z,\tau)$ be optimized variable
profit above. Assume an interior unique entry threshold $z_*(\tau)$ with
$V(z_*(\tau),\tau)=F$, positive investment there, and local smoothness
away from the financing kink. Entrants operate exactly for $z\geq z_*$.
The aggregate is
\[
 I(\tau)=\int_{z_*(\tau)}^{\overline z}i(z,b;\tau)f(z)\,dz,
\]
so Leibniz's rule gives
\[
 I'(0)=\int_{z_*(0)}^{\overline z}\dot i(z,b)f(z)\,dz
       -i(z_*(0),b;0)f(z_*(0))z_*'(0).
\]
The boundary term need not vanish even though the entering firm's net
profit is zero: its investment is positive.

For more detail, let $\eta\geq0$ be the multiplier on
$b-(1-\tau)i\geq0$. The Lagrangian is
$(z-1+\tau)i-qi^2/2+\eta[b-(1-\tau)i]$.
The envelope derivatives are $V_z=i$ and $V_\tau=(1+\eta)i$.
Implicit differentiation therefore yields
$z_*'=-(1+\eta_*)$ in this model. The entry contribution is
$i_*f(z_*)(1+\eta_*)\geq0$. This expression exposes a financing margin
that an incumbent-only average excludes. It relies on smooth threshold
crossing; discrete masses require finite differences instead.

\subsection*{General composition and equilibrium derivatives}
More generally assume
$\mu^\tau=\mu^0+\tau\nu+o(\tau)$ in total variation and
$i^\tau=i^0+\tau\dot i+o(\tau)$ uniformly, with finite measures and
bounded functions. Multiplying the two expansions proves
\[
 \frac{d}{d\tau}\int i^\tau\,d\mu^\tau\bigg|_0
 =\int\dot i\,d\mu^0+\int i^0\,d\nu.
\]
The signed measure $\nu$ need not integrate to zero because entry changes
the number of firms. This differs from the mass-conserving household
transition case. When endogenous states move continuously, total-variation
differentiability may be too strong; one can instead require weak
differentiability against the particular test functions used in the
aggregate, or integrate complete incumbent paths and add entry cohorts.
The explicit threshold example is directly covered by Leibniz's rule,
without pretending moving point masses have total-variation derivatives.

If prices $p(\tau)$ are determined by $F(p,\tau)=0$ with nonsingular
$F_p$, the policy derivative of the investment rule additionally contains
$i_p p'=-i_pF_p^{-1}F_\tau$. Taxes financing the subsidy must be included
in those equilibrium equations. The payment $\tau I$ is a transfer, but
its financing can alter household saving, wages and demand; it cannot be
ignored by invoking initial budget balance alone.

\subsection*{A defined financial-constraint decomposition}
Let $I(\tau,r)$ denote a fully re-cleared equilibrium, where $r=0$ retains
the financing restriction and $r=1$ removes it by an explicitly funded
intervention. Define the interaction
\[
 H=[I(\tau,1)-I(0,1)]-[I(\tau,0)-I(0,0)].
\]
This is the difference between two subsidy effects, not a regression
coefficient on baseline liquidity. Removing constraints before versus
after changing subsidy assigns the interaction to different channels.
For a symmetric two-intervention accounting, split it equally; equivalently
average each intervention's effect across its two possible orderings.
The two averaged components sum exactly to $I(\tau,1)-I(0,0)$ by
cancelling intermediate terms. This is an accounting definition, and each
of its four economies still requires identification.

\subsection*{Verification and remaining research}
The firm optimizer checks zero investment, slack finance and binding
finance; $\tau<1$ prevents a singular upfront-price constraint. The entry
formula includes newly operating firms, and financing is not counted as
an independent resource benefit. These results show exactly what a frozen-
distribution calculation misses. They do not identify a dynamic general-
equilibrium response from available empirical data or settle which fiscal
instrument gives the desired investment path. A next concrete step is to
estimate baseline-state investment derivatives and entry threshold shifts
jointly, then impose market-clearing restrictions with bounded price
responses. \textbf{Final status: UNRESOLVED.}

\subsection*{Source}
B. Moll, \emph{Lecture 11: Good Research Topics and Open Questions},
Distributional Macroeconomics lecture slides, PDF page 10 (firm cash flow,
investment and taxes).
\url{https://benjaminmoll.com/wp-content/uploads/2019/07/Lecture11_2149.pdf}.

\end{document}
